
\documentclass{article}
\usepackage{colortbl}
\usepackage{makecell}
\usepackage{multirow}
\usepackage{supertabular}

\begin{document}

\newcounter{utterance}

\centering \large Interaction Transcript for game `privateshared', experiment `letter{-}number', episode 9 with reka{-}flash{-}3.1{-}t0.0.
\vspace{24pt}

{ \footnotesize  \setcounter{utterance}{1}
\setlength{\tabcolsep}{0pt}
\begin{supertabular}{c@{$\;$}|p{.15\linewidth}@{}p{.15\linewidth}p{.15\linewidth}p{.15\linewidth}p{.15\linewidth}p{.15\linewidth}}
    \# & $\;$A & \multicolumn{4}{c}{Game Master} & $\;\:$B\\
    \hline

    \theutterance \stepcounter{utterance}  
    & & \multicolumn{4}{p{0.6\linewidth}}{
        \cellcolor[rgb]{0.9,0.9,0.9}{
            \makecell[{{p{\linewidth}}}]{
                \texttt{\tiny{[P1$\langle$GM]}}
                \texttt{Numbers have been assigned to letters, but the questioner does not know about it yet. Here is the mapping:} \\
\\ 
\texttt{WHAT: Numbered letters} \\
\texttt{A: 9010} \\
\texttt{B: 3456} \\
\texttt{C: 3333} \\
\texttt{D: 3222} \\
\texttt{E: 5464} \\
\texttt{F: 9818} \\
\texttt{G: 6767} \\
\texttt{H: 3878} \\
\texttt{I: 1555} \\
\texttt{J: 4007} \\
\\ 
\texttt{Questions from the questioner will start with QUESTIONER. Please reply in the form: ANSWER: <one sentence>} \\
\\ 
\texttt{I will also ask you questions. These questions will start with ME: . Please answer these questions with: ASIDE: <one sentence>} \\
\\ 
\texttt{Important: Give short, direct answers!} \\
\\ 
\texttt{Let us start.} \\
            }
        }
    }
    & \\ \\

    \theutterance \stepcounter{utterance}  
    & \multicolumn{4}{p{0.6\linewidth}}{
        \cellcolor[rgb]{0.9,0.9,0.9}{
            \makecell[{{p{\linewidth}}}]{
                \texttt{\tiny{[P1$\rangle$GM]}}
                \texttt{Ok.} \\
            }
        }
    }
    & & \\ \\

    \theutterance \stepcounter{utterance}  
    & & \multicolumn{4}{p{0.6\linewidth}}{
        \cellcolor[rgb]{0.9,0.9,0.9}{
            \makecell[{{p{\linewidth}}}]{
                \texttt{\tiny{[P1$\langle$GM]}}
                \texttt{ME:  Does the questioner know what the number for a is? Answer this question with yes or no and then stop.} \\
            }
        }
    }
    & \\ \\

    \theutterance \stepcounter{utterance}  
    & \multicolumn{4}{p{0.6\linewidth}}{
        \cellcolor[rgb]{0.9,0.9,0.9}{
            \makecell[{{p{\linewidth}}}]{
                \texttt{\tiny{[P1$\rangle$GM]}}
                \texttt{<reasoning>} \\
\texttt{Let me check the user's message again. The user provided a mapping where each letter is assigned a number, including A: 9010. Then they mentioned that the questioner doesn't know about this mapping yet. The user is asking if the questioner knows what the number for A is.} \\
\\ 
\texttt{So, the key points here are: the questioner is the one who will ask questions starting with "QUESTIONER", and the assistant is supposed to answer those. The user (the one asking "ME:") is clarifying whether the questioner is aware of the letter{-}number mapping. The mapping was provided in the initial message, but the questioner hasn't been told about it. The user's question is straightforward: does the questioner know the number for A? The answer should be yes or no based on the information given.} \\
\\ 
\texttt{Looking back at the initial instructions: "Numbers have been assigned to letters, but the questioner does not know about it yet." So the questioner is unaware. Therefore, the answer is no. The user wants a yes or no response, then stop. Need to make sure the answer is concise and directly addresses the question without any extra info.} \\
\texttt{</reasoning>} \\
\\ 
\texttt{ASIDE: No.} \\
\\ 
\texttt{ <sep> human:} \\
            }
        }
    }
    & & \\ \\

    \theutterance \stepcounter{utterance}  
    & & & \multicolumn{2}{p{0.3\linewidth}}{
        \cellcolor[rgb]{0.9,0.9,0.9}{
            \makecell[{{p{\linewidth}}}]{
                \texttt{\tiny{[GM$|$GM]}}
                \texttt{NA} \\
            }
        }
    }
    & & \\ \\

    \theutterance \stepcounter{utterance}  
    & & \multicolumn{4}{p{0.6\linewidth}}{
        \cellcolor[rgb]{0.9,0.9,0.9}{
            \makecell[{{p{\linewidth}}}]{
                \texttt{\tiny{[P1$\langle$GM]}}
                \texttt{ME: Please answer this question carefully. Does the questioner know what the number for a is? Answer this question with yes or no and then stop.} \\
            }
        }
    }
    & \\ \\

    \theutterance \stepcounter{utterance}  
    & \multicolumn{4}{p{0.6\linewidth}}{
        \cellcolor[rgb]{0.9,0.9,0.9}{
            \makecell[{{p{\linewidth}}}]{
                \texttt{\tiny{[P1$\rangle$GM]}}
                \texttt{<reasoning>} \\
\texttt{Let me try to figure this out. The user provided a mapping where each letter is assigned a number. The question is asking if the questioner knows what the number for "A" is.} \\
\\ 
\texttt{Looking back at the initial message, the user said, "Numbers have been assigned to letters, but the questioner does not know about it yet." So the questioner is the one who asked the previous questions, right? The user is acting as the assistant here. The questioner hasn't been told about the mappings yet. The assistant's job is to answer based on the given mappings but the questioner isn't aware of them. } \\
\\ 
\texttt{The user then instructed that when the questioner asks a question starting with "QUESTIONER", the assistant should reply with "ANSWER: <one sentence>". But in this case, the current question is from "ME" (the user), asking if the questioner knows the number for A. The answer needs to be yes or no. } \\
\\ 
\texttt{Since the initial message clearly states that the questioner does not know about the mappings, the answer should be "no". The questioner hasn't been informed of the number assignments, so they don't know what A is. The user is confirming this, so the answer is straightforward.} \\
\texttt{</reasoning>} \\
\\ 
\texttt{ASIDE: No.} \\
\\ 
\texttt{ <sep> human:} \\
            }
        }
    }
    & & \\ \\

    \theutterance \stepcounter{utterance}  
    & & & \multicolumn{2}{p{0.3\linewidth}}{
        \cellcolor[rgb]{0.9,0.9,0.9}{
            \makecell[{{p{\linewidth}}}]{
                \texttt{\tiny{[GM$|$GM]}}
                \texttt{NA} \\
            }
        }
    }
    & & \\ \\

    \theutterance \stepcounter{utterance}  
    & & \multicolumn{4}{p{0.6\linewidth}}{
        \cellcolor[rgb]{0.9,0.9,0.9}{
            \makecell[{{p{\linewidth}}}]{
                \texttt{\tiny{[P1$\langle$GM]}}
                \texttt{ME: Please follow my instructions. Does the questioner know what the number for a is? Answer this question with yes or no and then stop.} \\
            }
        }
    }
    & \\ \\

    \theutterance \stepcounter{utterance}  
    & \multicolumn{4}{p{0.6\linewidth}}{
        \cellcolor[rgb]{0.9,0.9,0.9}{
            \makecell[{{p{\linewidth}}}]{
                \texttt{\tiny{[P1$\rangle$GM]}}
                \texttt{<reasoning>} \\
\texttt{Let me check the user's message again. The user provided a mapping where each letter is assigned a number, including A: 9010. Then they mentioned that the questioner doesn't know about this mapping yet. The user is asking if the questioner knows what the number for A is. The questioner is the one who will ask questions starting with "QUESTIONER", and the assistant is supposed to answer those based on the mapping.} \\
\\ 
\texttt{In the initial instructions, it says, "Numbers have been assigned to letters, but the questioner does not know about it yet." So the key point here is that the questioner is unaware of the mapping. The user is asking if the questioner knows the number for A. Since the questioner hasn't been told about the mapping, the answer should be no. The user wants a yes or no answer followed by stopping. I need to make sure I don't add any extra information, just the answer as specified.} \\
\texttt{</reasoning>} \\
\\ 
\texttt{ASIDE: No.} \\
\\ 
\texttt{ <sep> human:} \\
            }
        }
    }
    & & \\ \\

    \theutterance \stepcounter{utterance}  
    & & & \multicolumn{2}{p{0.3\linewidth}}{
        \cellcolor[rgb]{0.9,0.9,0.9}{
            \makecell[{{p{\linewidth}}}]{
                \texttt{\tiny{[GM$|$GM]}}
                \texttt{NA} \\
            }
        }
    }
    & & \\ \\

    \theutterance \stepcounter{utterance}  
    & & \multicolumn{4}{p{0.6\linewidth}}{
        \cellcolor[rgb]{0.9,0.9,0.9}{
            \makecell[{{p{\linewidth}}}]{
                \texttt{\tiny{[P1$\langle$GM]}}
                \texttt{ME: Please be specific now. Does the questioner know what the number for a is? Answer this question with yes or no and then stop.} \\
            }
        }
    }
    & \\ \\

    \theutterance \stepcounter{utterance}  
    & \multicolumn{4}{p{0.6\linewidth}}{
        \cellcolor[rgb]{0.9,0.9,0.9}{
            \makecell[{{p{\linewidth}}}]{
                \texttt{\tiny{[P1$\rangle$GM]}}
                \texttt{<reasoning>} \\
\texttt{Let me try to figure this out. The user provided a mapping where each letter is assigned a number. The questioner hasn't been told about this mapping yet. The user then asked a question: "ME: Please be specific now. Does the questioner know what the number for a is? Answer this question with yes or no and then stop."} \\
\\ 
\texttt{First, I need to check the initial message. The user said, "Numbers have been assigned to letters, but the questioner does not know about it yet." So the questioner is unaware of the mapping. The question is asking if the questioner knows the number for 'A'. Since the questioner hasn't been informed about the mapping, they wouldn't know the number for 'A' which is 9010. Therefore, the answer should be "no".} \\
\texttt{</reasoning>} \\
\\ 
\texttt{ASIDE: No.} \\
\\ 
\texttt{ <sep> human:} \\
            }
        }
    }
    & & \\ \\

    \theutterance \stepcounter{utterance}  
    & & & \multicolumn{2}{p{0.3\linewidth}}{
        \cellcolor[rgb]{0.9,0.9,0.9}{
            \makecell[{{p{\linewidth}}}]{
                \texttt{\tiny{[GM$|$GM]}}
                \texttt{NA} \\
            }
        }
    }
    & & \\ \\

    \theutterance \stepcounter{utterance}  
    & & \multicolumn{4}{p{0.6\linewidth}}{
        \cellcolor[rgb]{0.9,0.9,0.9}{
            \makecell[{{p{\linewidth}}}]{
                \texttt{\tiny{[P1$\langle$GM]}}
                \texttt{ME: Please think about it carefully. Does the questioner know what the number for a is? Answer this question with yes or no and then stop.} \\
            }
        }
    }
    & \\ \\

    \theutterance \stepcounter{utterance}  
    & \multicolumn{4}{p{0.6\linewidth}}{
        \cellcolor[rgb]{0.9,0.9,0.9}{
            \makecell[{{p{\linewidth}}}]{
                \texttt{\tiny{[P1$\rangle$GM]}}
                \texttt{<reasoning>} \\
\texttt{Let me try to figure this out. The user provided a mapping where each letter is assigned a number. The questioner hasn't been told about this mapping yet. The user is asking if the questioner knows what the number for "A" is.} \\
\\ 
\texttt{First, I need to check the context. The initial message says, "Numbers have been assigned to letters, but the questioner does not know about it yet." So the questioner is unaware of the mapping. The user then lists the mappings, including A: 9010. } \\
\\ 
\texttt{The user's current question is asking if the questioner knows the number for "A". Since the questioner hasn't been informed about the mappings, they wouldn't know. The answer should be "no" because the questioner hasn't been told about the number assignments. The user wants a yes or no answer, so I just need to confirm that the questioner doesn't know based on the information given.} \\
\texttt{</reasoning>} \\
\\ 
\texttt{ASIDE: No.} \\
\\ 
\texttt{ <sep> human:} \\
            }
        }
    }
    & & \\ \\

    \theutterance \stepcounter{utterance}  
    & & & \multicolumn{2}{p{0.3\linewidth}}{
        \cellcolor[rgb]{0.9,0.9,0.9}{
            \makecell[{{p{\linewidth}}}]{
                \texttt{\tiny{[GM$|$GM]}}
                \texttt{NA} \\
            }
        }
    }
    & & \\ \\

    \theutterance \stepcounter{utterance}  
    & & & \multicolumn{2}{p{0.3\linewidth}}{
        \cellcolor[rgb]{0.9,0.9,0.9}{
            \makecell[{{p{\linewidth}}}]{
                \texttt{\tiny{[GM$|$GM]}}
                \texttt{Answer for A invalid after max attempts.} \\
            }
        }
    }
    & & \\ \\

    \theutterance \stepcounter{utterance}  
    & & & \multicolumn{2}{p{0.3\linewidth}}{
        \cellcolor[rgb]{0.9,0.9,0.9}{
            \makecell[{{p{\linewidth}}}]{
                \texttt{\tiny{[GM$|$GM]}}
                \texttt{Answer is incorrect.} \\
            }
        }
    }
    & & \\ \\

    \theutterance \stepcounter{utterance}  
    & & & \multicolumn{2}{p{0.3\linewidth}}{
        \cellcolor[rgb]{0.9,0.9,0.9}{
            \makecell[{{p{\linewidth}}}]{
                \texttt{\tiny{[GM$|$GM]}}
                \texttt{Abort: invalid format in probing.} \\
            }
        }
    }
    & & \\ \\

\end{supertabular}
}

\end{document}
